\begin{lemma}[Certified restriction to every earlier positive rung]
Let $G$ be a finite loopless multigraph with palette $[k]$ and
distinguished edge $e_0$.  Let $\varphi_{\mathrm{in}}$ be an input
partial edge coloring, and let $\mathcal E$ be an
\noderef{ExtendedTashkinovTree} on these data, with displayed terminal
tree $T$, displayed generating coloring $\varphi_n$, rung number $n$,
certified \noderef{TashkinovSeriesData} $\mathcal D$, and
\noderef{MaximumProperty}.
For every integer $r$ satisfying
\[
                         1\leq r<n,
\]
there is an actual extended Tashkinov tree $\mathcal E^{\leq r}$ on the
same graph, palette, distinguished edge, and input coloring whose
displayed terminal tree is exactly the original tree slot
$\mathcal D.\mathit{tree}(r+1)$ and whose displayed generating coloring
is exactly $\mathcal D.\mathit{coloring}(r)$.  Writing
$\mathcal D^{\leq r}$ for its certified series data, it has all of the
following properties.
\begin{enumerate}
\item Its rung number is exactly $r$, and it has maximum property.
\item Under the canonical order-preserving inclusion
$\iota_T:\operatorname{Fin}(r+2)\to\operatorname{Fin}(n+2)$, every
retained tree slot agrees pointwise:
\[
 \mathcal D^{\leq r}.\mathit{tree}(a)
   =\mathcal D.\mathit{tree}(\iota_T(a)).
\]
\item Under the canonical inclusion
$\iota_D:\operatorname{Fin}(r+1)\to\operatorname{Fin}(n+1)$, the
coloring, raw stored-color set, exact connecting-edge set, extension
kind, chosen vertex, outside vertex, delta, and gamma arrays all agree
pointwise with their original arrays.
\item At every retained data index $a$, the computed stored-color set is
preserved exactly:
\[
 \noderef{TashkinovStoredColors}(\mathcal D^{\leq r},a)
  =\noderef{TashkinovStoredColors}(\mathcal D,\iota_D(a)).
\]
\item At every retained data index, the
\noderef{TashkinovSeriesStep} certificate is the original certificate at
the included index after dependent transport along these equalities; in
particular, the two certificate terms are heterogeneously equal.
\end{enumerate}
No restriction ETT is asserted at target rung zero.
\end{lemma}
\begin{proof}
We prove the assertion simultaneously for all $\mathcal E$ and all
admissible $r$ by strong induction on the positive gap
\[
                              d=n-r.
\]
Because $r\geq1$ and $r<n$, one has $n>1$, so
\noderef{CertifiedETTPredecessor} applies to $\mathcal E$.  It returns an
actual maximum ETT $\mathcal P$ whose rung number is $n-1$, whose
displayed terminal tree and generating coloring are exactly the original
slots $\mathcal D.\mathit{tree}(n)$ and
$\mathcal D.\mathit{coloring}(n-1)$, and whose retained series data,
computed stored-color sets, and dependent step certificates agree with
those of $\mathcal E$ under the canonical inclusions.  Denote its series
data by $\mathcal D^-$.  The cited result also supplies the exact rung
equation and maximum property for $\mathcal P$.

Suppose first that $d=1$.  Natural-number arithmetic gives $n=r+1$ and
$n-1=r$.  Thus $\mathcal P$ already has the required rung, displayed
tree $\mathcal D.\mathit{tree}(r+1)$, and displayed coloring
$\mathcal D.\mathit{coloring}(r)$.  The tree inclusion
$\operatorname{Fin}((n-1)+2)\to\operatorname{Fin}(n+2)$ and the data
inclusion
$\operatorname{Fin}((n-1)+1)\to\operatorname{Fin}(n+1)$ supplied by
\noderef{CertifiedETTPredecessor} become exactly $\iota_T$ and
$\iota_D$ after substituting $n-1=r$.  Its pointwise equalities therefore
give items 2--4, its heterogeneous certificate equality gives item 5,
and its rung equation and inherited maximum property give item 1.  This
proves the base case without constructing a zero-rung ETT.

Now suppose $d>1$.  Then $r<n-1$, while $r\geq1$, so the induction
hypothesis applies to the maximum ETT $\mathcal P$ at the same target
rung $r$.  Its gap is
\[
                    (n-1)-r=(n-r)-1=d-1<d.
\]
It gives an ETT $\mathcal R$ of rung $r$, initially displayed on
$\mathcal D^-.\mathit{tree}(r+1)$ and
$\mathcal D^-.\mathit{coloring}(r)$, with maximum property and the full
five-part agreement package relative to $\mathcal D^-$.

The predecessor tree equality at index $r+1$ and coloring equality at
index $r$ identify those two displayed objects with
$\mathcal D.\mathit{tree}(r+1)$ and
$\mathcal D.\mathit{coloring}(r)$.  Transport $\mathcal R$ along these
two exact equalities in the dependent displayed-tree and
displayed-coloring parameters of
\noderef{ExtendedTashkinovTree}.  Call the transported ETT
$\mathcal E^{\leq r}$.  Equality elimination changes no rung or retained
series datum; it transports the maximum-property proof and the recursive
agreement proofs to their required types.  Hence
$\mathcal E^{\leq r}$ is an actual ETT on the original input data, has
rung $r$, has the two exact displayed objects asserted in the statement,
and has maximum property.

It remains to verify that the two layers of retained-data equalities have
the single canonical inclusions stated in the conclusion.  Let
\[
 \jmath_T:\operatorname{Fin}(r+2)\to
              \operatorname{Fin}((n-1)+2),\qquad
 \kappa_T:\operatorname{Fin}((n-1)+2)\to
              \operatorname{Fin}(n+2)
\]
be respectively the tree inclusions from the recursive restriction and
the predecessor restriction.  Both are the `Fin.castLE' map: they retain
the numerical value of their argument and merely enlarge its upper
bound.  Consequently, for every $a$, the elements
$\kappa_T(\jmath_T(a))$ and $\iota_T(a)$ have the same value.  Extensionality
for finite indices gives
\[
                         \kappa_T(\jmath_T(a))=\iota_T(a).       \tag{1}
\]
The corresponding data inclusions
\[
 \jmath_D:\operatorname{Fin}(r+1)\to
              \operatorname{Fin}((n-1)+1),\qquad
 \kappa_D:\operatorname{Fin}((n-1)+1)\to
              \operatorname{Fin}(n+1)
\]
also preserve values, and the same argument gives
\[
                         \kappa_D(\jmath_D(a))=\iota_D(a).       \tag{2}
\]

For a tree index $a$, the recursive equality identifies the retained
tree with $\mathcal D^-.\mathit{tree}(\jmath_T(a))$; the predecessor
equality identifies this with
$\mathcal D.\mathit{tree}(\kappa_T(\jmath_T(a)))$; and (1) replaces the
composite index by $\iota_T(a)$.  This proves item 2.  Apply the identical
three-equality chain using (2) separately to the coloring, raw stored
set, exact connecting-edge set, kind, chosen vertex, outside vertex,
delta, and gamma entries.  This proves every clause of item 3, including
equality of individual parallel-edge identities.

For the computed set at $a$, the induction hypothesis gives equality
with
$\noderef{TashkinovStoredColors}(\mathcal D^-,\jmath_D(a))$.
The predecessor theorem gives equality of that set with the original
computed set at $\kappa_D(\jmath_D(a))$, and (2) changes that index to
$\iota_D(a)$.  Transitivity proves item 4; no equality of raw stored slots
is being substituted for equality of the computed bounded union and
missing-color deletion.

Finally, the induction hypothesis gives heterogeneous equality between
the retained certificate and the certificate of $\mathcal P$ at
$\jmath_D(a)$.  The predecessor theorem gives heterogeneous equality
from that certificate to the original certificate at
$\kappa_D(\jmath_D(a))$.  Heterogeneous equality is transitive, and
transport along (2) identifies the latter dependent term with the
original certificate at $\iota_D(a)$.  Equality transport in the two
displayed parameters does not change this proposition-valued
certificate; proof irrelevance identifies any transported proof terms.
This proves item 5.

Thus all five conclusions hold for gap $d$.  Strong induction proves them
for every $1\leq r<n$, completing the construction.
\end{proof}
